\section{Robustness of the rescue flip}
\label{sec:robustness}

Sections~\ref{sec:doubleedge} through~\ref{sec:reachingedge} have built and tested the rescue flip at one specific parameter combination ($\Nzero = 1000$, $l = 100$, $\mu = 10^{-4}$). A reader who has accepted the race mechanism in Section~\ref{sec:reachingedge} would reasonably ask whether the rescue flip is a robust qualitative feature of the model, or an artifact of this one setting. We sweep the three parameters held fixed in those sections (edge width $l$, core activity $\bCore$, and mutation rate $\mu$) and ask two questions of each sweep: does the flip persist across the parameter range, and does the crossover dose shift in ways the race mechanism from Section~\ref{sec:reachingedge} would predict?

Figure~\ref{fig:sensL} sweeps the edge width $l$ across $\{25, 50, 100, 200, 400\}$, holding all other parameters at their defaults. Panel A shows the dose curves for the biofilm at each $l$ alongside the well-mixed reference; Panel B shows the corresponding $\Delta P(\mathrm{rescue}) = P_{\mathrm{BF}} - P_{\mathrm{WM}}$. The flip persists across a sixteen-fold change in edge width, but its magnitude scales strongly with $l$. At smaller $l$ the flip is sharper, with the biofilm's high-dose advantage growing as $l$ shrinks: at $\dEdge = 2.8$, $\Delta P \approx 0.013$ for $l = 25$ but is essentially zero at $l = 400$, where the biofilm and well-mixed populations are indistinguishable across the entire dose range. The crossover dose itself shifts only modestly, from roughly $\dEdge \approx 1.5$ at $l = 200$ down to $\dEdge \approx 1.4$ at $l = 25$.

\begin{figure}[h!]
\centering
\includegraphics[width=0.95\textwidth]{sensL_curves.png}
\caption{Sensitivity of the rescue flip to edge width $l$. Panel A: $P(\mathrm{rescue})$ vs dose for the biofilm at five values of $l \in \{25, 50, 100, 200, 400\}$ and the well-mixed reference (open circles, dashed). Panel B: difference relative to the well-mixed reference, $\Delta P(\mathrm{rescue}) = P_{\mathrm{BF}} - P_{\mathrm{WM}}$, with the dotted horizontal line at zero marking the no-difference baseline. Each point averages $20{,}000$ independent simulations per parameter combination. \emph{Parameters}: $\Nzero = 1000$, $\bCore = 0.2$, $\mu = 10^{-4}$.}
\label{fig:sensL}
\end{figure}

Both shifts can be read through the race mechanism from Section~\ref{sec:reachingedge}. Smaller $l$ has two competing effects on the race. The boundary speed is $\sE \cdot l$, so smaller $l$ slows boundary retreat in absolute cell terms, which by itself would reduce the fraction of core mutations that reach the edge. But smaller $l$ also makes the core larger relative to the total population (it is $\Nzero - l$ cells), which means more core births per unit time and a longer Phase 1 over which to produce them: the duration $T_1 = (\Nzero - l) / (\sE \cdot l)$ increases as $l$ decreases. The larger pool of core mutations more than compensates for the slower boundary retreat. The net effect tilts the race toward the biofilm at smaller $l$, which is what Panel B shows.

Figure~\ref{fig:sensCore} sweeps the core metabolic activity $\bCore$ across $\{0, 0.05, 0.1, 0.2, 0.4, 0.8\}$, holding all other parameters at their defaults. Panel A shows the dose curves for the biofilm at each $\bCore$ alongside the well-mixed reference; Panel B shows the corresponding $\Delta P(\mathrm{rescue})$. At $\bCore = 0$ (dormant core), the biofilm stays strictly below the well-mixed reference at every dose; there is no crossover. At every $\bCore > 0$ in the sweep, the biofilm crosses the well-mixed reference at some dose and stays at or above it at higher doses. The crossover dose shifts predictably with $\bCore$: from $\dEdge \approx 1.6$ at $\bCore = 0.05$ down to $\dEdge \approx 1.2$ at $\bCore = 0.8$. The high-dose advantage also grows with $\bCore$: at $\dEdge = 2.8$, $\Delta P$ at $\bCore = 0.8$ is about four times its value at $\bCore = 0.05$.

\begin{figure}[h!]
\centering
\includegraphics[width=0.95\textwidth]{sensCore_curves.png}
\caption{Sensitivity of the rescue flip to core metabolic activity $\bCore$. Panel A: $P(\mathrm{rescue})$ vs dose for the biofilm at six values of $\bCore \in \{0, 0.05, 0.1, 0.2, 0.4, 0.8\}$ and the well-mixed reference (open circles, dashed). Panel B: difference relative to the well-mixed reference, $\Delta P(\mathrm{rescue}) = P_{\mathrm{BF}} - P_{\mathrm{WM}}$, with the dotted horizontal line at zero marking the no-difference baseline. Each point averages $20{,}000$ independent simulations per parameter combination. \emph{Parameters}: $\Nzero = 1000$, $l = 100$, $\mu = 10^{-4}$.}
\label{fig:sensCore}
\end{figure}

The shifts can be read through the race mechanism from Section~\ref{sec:reachingedge}. The dormant case ($\bCore = 0$) is the cleanest test: with no core births, no R lineages are ever produced in the core for the boundary to catch, and the biofilm gains no advantage at any dose. The $\bCore > 0$ cases trade off two effects of increasing the core's metabolic rate. Higher $\bCore$ produces more core mutations per unit time, which feeds more lineages into the race. But higher $\bCore$ also accelerates drift in the core, so a smaller fraction of those lineages survive long enough for the boundary to reach them (Figure~\ref{fig:reaching} showed this directly: at $\bCore = 0.8$, fewer core lineages reach the edge as a fraction than at $\bCore = 0.2$). The absolute mutation count rises faster than the fraction lost, so the total flow of core lineages arriving at the edge increases with $\bCore$, and the biofilm's high-dose advantage grows accordingly.

Figure~\ref{fig:sensMu} sweeps the mutation rate $\mu$ across $\{10^{-5}, 10^{-4}, 10^{-3}\}$, holding all other parameters at their defaults. Mutation rate scales the absolute rescue probability strongly: at $\mu = 10^{-3}$ both populations rescue close to a substantial fraction of the time at low dose, while at $\mu = 10^{-5}$ rescue is rare for both at every dose. The qualitative flip persists at $\mu = 10^{-4}$ (our default, with crossover at $\dEdge \approx 1.4$) and at $\mu = 10^{-3}$ (crossover at $\dEdge \approx 1.5$, with the high-dose advantage about an order of magnitude larger than at the default $\mu$). At $\mu = 10^{-5}$, both populations rescue so rarely that the absolute $\Delta P$ stays near zero across the sweep; resolving the flip cleanly at this mutation rate would require many more replicates than we have. Within the range where the model produces measurable rescue events, the qualitative shape of the flip is robust to $\mu$: the mutation rate sets how often any rescue happens, but does not change which population gets the larger flow of core-born lineages arriving at the edge at high dose.

\begin{figure}[h!]
\centering
\includegraphics[width=0.95\textwidth]{sensMu_curves.png}
\caption{Sensitivity of the rescue flip to mutation rate $\mu$. Panel A: $P(\mathrm{rescue})$ vs dose for the biofilm at three values of $\mu \in \{10^{-5}, 10^{-4}, 10^{-3}\}$ and the well-mixed reference at $\mu = 10^{-4}$ (open circles, dashed). Panel B: difference relative to the well-mixed reference at each $\mu$, $\Delta P(\mathrm{rescue}) = P_{\mathrm{BF}} - P_{\mathrm{WM}}$, with the dotted horizontal line at zero. Each point averages $20{,}000$ independent simulations per parameter combination. \emph{Parameters}: $\Nzero = 1000$, $l = 100$, $\bCore = 0.2$.}
\label{fig:sensMu}
\end{figure}

Across all three sweeps, the rescue flip is a robust qualitative feature of the model. The flip persists at every edge width $l$ we tested (a sixteen-fold change in $l$), at every nonzero core activity $\bCore$, and at every mutation rate $\mu$ for which the model produces measurable rescue events. The locations of the flip (the crossover doses) shift in directions that the race mechanism from Section~\ref{sec:reachingedge} predicts: changes in $l$ alter both the boundary speed and the duration of Phase 1, changes in $\bCore$ alter both how many core mutations are produced and how fast they drift to extinction, and changes in $\mu$ scale the overall mutation supply without changing the race dynamics. The biofilm wins where these factors combine to produce a flow of core lineages arriving at the edge large enough to overcome the biofilm's per-mutation establishment penalty.

Sections~\ref{sec:doubleedge} through~\ref{sec:robustness} have now built the rescue flip, measured its mechanism, and tested its robustness across the three parameters that govern it. What remains is the question Section~\ref{sec:intro} actually opened with: what does this tell us about treating biofilm infections, and in particular about behavior-modifying interventions that disrupt biofilm spatial structure? The next section returns to that motivating question with the mechanism in hand.
